% L-GTA Architecture Diagram
% Add the following to your main document preamble:
%   \usepackage{tikz}
%   \usetikzlibrary{positioning, calc, arrows.meta, matrix, fit}
% Compile with pdflatex or include in a main document via % L-GTA Architecture Diagram
% Add the following to your main document preamble:
%   \usepackage{tikz}
%   \usetikzlibrary{positioning, calc, arrows.meta, matrix, fit}
% Compile with pdflatex or include in a main document via % L-GTA Architecture Diagram
% Add the following to your main document preamble:
%   \usepackage{tikz}
%   \usetikzlibrary{positioning, calc, arrows.meta, matrix, fit}
% Compile with pdflatex or include in a main document via % L-GTA Architecture Diagram
% Add the following to your main document preamble:
%   \usepackage{tikz}
%   \usetikzlibrary{positioning, calc, arrows.meta, matrix, fit}
% Compile with pdflatex or include in a main document via \input{docs/lgta_architecture}

\begin{figure*}[t]
\centering
\resizebox{\textwidth}{!}{%
\begin{tikzpicture}[
    >=Stealth,
    % ── colour palette ──
    encoder_col/.style  = {fill=blue!6,   draw=blue!45!black,  line width=0.8pt},
    decoder_col/.style  = {fill=orange!6, draw=orange!50!black, line width=0.8pt},
    latent_col/.style   = {fill=green!6,  draw=green!45!black,  line width=0.8pt},
    transf_col/.style   = {fill=violet!6, draw=violet!45!black, line width=0.8pt},
    io_col/.style       = {fill=gray!6,   draw=gray!55,         line width=0.8pt},
    loss_col/.style     = {fill=red!5,    draw=red!45!black,    line width=0.8pt},
    equiv_col/.style    = {fill=teal!5,   draw=teal!50!black,   line width=0.8pt,
                           dashed, dash pattern=on 3pt off 2pt},
    ctx_col/.style      = {fill=yellow!6, draw=yellow!50!black, line width=0.8pt},
    % ── shapes ──
    bigblock/.style = {
        rectangle, rounded corners=4pt,
        minimum height=3.4cm, minimum width=4.2cm,
        align=center},
    medblock/.style = {
        rectangle, rounded corners=4pt,
        minimum height=3.4cm, minimum width=3.6cm,
        align=center},
    innerblock/.style = {
        rectangle, rounded corners=2pt,
        minimum height=0.6cm, minimum width=3.0cm,
        align=center, font=\footnotesize, line width=0.5pt},
    ioblock/.style = {
        rectangle, rounded corners=3pt,
        minimum height=1.2cm, minimum width=1.8cm,
        align=center, font=\normalsize, io_col},
    prepblock/.style = {
        rectangle, rounded corners=3pt,
        minimum height=1.2cm, minimum width=2.4cm,
        align=center, font=\footnotesize, io_col},
    arrow/.style = {->, line width=1.2pt, draw=gray!60},
    thinarrow/.style = {->, line width=0.8pt, draw=gray!50},
    dasharrow/.style = {->, line width=0.8pt, draw=teal!50!black,
                        dashed, dash pattern=on 3pt off 2pt},
    dimlabel/.style = {font=\footnotesize\itshape, text=gray!65!black},
    seclabel/.style = {font=\normalsize\bfseries},
]

% =====================================================================
%  ROW 1 (top, left to right): Input → Encoder → Latent
% =====================================================================

\node[ioblock] (input) at (0, 0)
    {Input\\[2pt]$z \in \mathbb{R}^{T \times S}$};

\node[bigblock, encoder_col] (enc) at (5.4, 0) {};
\node[seclabel, text=blue!45!black, anchor=north]
    at ([yshift=-4pt]enc.north) {Encoder\; $\mathcal{E}$};
\node[innerblock, encoder_col] (enc_bilstm)
    at ([yshift=0.05cm]enc.center)
    {Bi-LSTM $\to$ LayerNorm};
\node[innerblock, encoder_col, below=0.3cm of enc_bilstm] (enc_attn)
    {Temporal Self-Attention\\[-1pt]
     {\scriptsize pos.\ embed \textbullet{} MHA \textbullet{} FFN}};

\node[medblock, latent_col, minimum height=3.8cm] (lat) at (10.4, 0) {};
\node[seclabel, text=green!38!black, anchor=north]
    at ([yshift=-4pt]lat.north) {Latent Space};
\node[innerblock, latent_col, minimum width=2.4cm]
    (lat_mu) at ([yshift=0.1cm]lat.center)
    {$\mu_t,\;\log\Sigma_t$};
\node[innerblock, latent_col, minimum width=2.4cm,
      below=0.25cm of lat_mu] (lat_sample)
    {Sampling\\[-1pt]{\scriptsize reparam.\ trick}};
\node[dimlabel, below=0.1cm of lat_sample]
    {$v_t \sim \mathcal{N}(\mu_t,\, \Sigma_t)$};

% =====================================================================
%  TURN: Latent Space → Transforms (going down)
% =====================================================================

\node[medblock, transf_col, minimum height=4.2cm, minimum width=3.9cm] (transf) at (10.4, -5.2) {};
\node[seclabel, text=violet!45!black, anchor=north]
    at ([yshift=-4pt]transf.north) {Latent Transforms};
\node[innerblock, transf_col, minimum width=2.4cm]
    (transf_fn) at ([yshift=0.45cm]transf.center)
    {$v' = \mathcal{T}_\eta(v,\;\eta)$};
\node[font=\footnotesize, text=violet!40!black, align=center,
      below=0.25cm of transf_fn] (transf_list)
    {jitter \textbullet{} scaling \textbullet{} mag.\ warp\\
     drift \textbullet{} trend};
\node[dimlabel, align=center, below=0.25cm of transf_list]
    {chainable: $\mathcal{T}_n \!\circ\! \cdots \!\circ\! \mathcal{T}_1$};

% =====================================================================
%  ROW 2 (bottom, right to left): Decoder → Output
% =====================================================================

\node[bigblock, decoder_col, minimum height=4.2cm] (dec) at (5.4, -5.2) {};
\node[seclabel, text=orange!50!black, anchor=north]
    at ([yshift=-4pt]dec.north) {Decoder\; $\mathcal{D}$};
\node[font=\footnotesize\itshape, text=orange!45!black,
      anchor=north] at ([yshift=-0.5cm]dec.north)
    {$+\; z_{\mathrm{proj}}(v')$ skip connection};
\node[innerblock, decoder_col] (dec_bilstm)
    at ([yshift=-0.05cm]dec.center) {Bi-LSTM};
\node[innerblock, decoder_col, below=0.3cm of dec_bilstm] (dec_fc)
    {FC };

\node[ioblock] (output) at (0, -5.2)
    {Output\\[2pt]$\tilde{z} \in \mathbb{R}^{T \times S}$};

% =====================================================================
%  CONTEXT (aligned with encoder and decoder)
% =====================================================================
\node[prepblock, ctx_col, minimum width=2.8cm, minimum height=1.2cm]
    (ctx) at (5.4, -2.4)
    {Context $c$\\[1pt]{\scriptsize trigonometric time features}};

% =====================================================================
%  MAIN FLOW ARROWS
% =====================================================================

% Top row (left → right)
\draw[arrow] (input) -- (enc);
\draw[arrow] (enc)   -- (lat);

% Encoder internal
\draw[thinarrow] (enc_bilstm) -- (enc_attn);

% Turn down (latent → transforms)
\draw[arrow] (lat.south) -- (transf.north);

% Bottom row (right → left)
\draw[arrow] (transf.west) -- (dec.east);
\draw[arrow] (dec) -- (output);

% =====================================================================
%  CONTEXT ARROWS: top of context → encoder bottom, bottom of context → decoder top
% =====================================================================
\draw[->, line width=1.0pt, draw=yellow!55!black]
    (ctx.north) -- (enc.south);

\draw[->, line width=1.0pt, draw=yellow!55!black]
    (ctx.south) -- (dec.north);

% =====================================================================
%  TRAINING LOSS
% =====================================================================
\node[rectangle, rounded corners=4pt, loss_col,
      minimum height=1.1cm,
      align=center, font=\small, inner sep=8pt]
    (loss) at (4.8, -8.5)
    {$\mathcal{L}
       = \mathcal{L}_{\mathrm{recon}}
       + \beta\,\mathcal{L}_{\mathrm{KL}}
       + \lambda_{\mathrm{equiv}}\,\mathcal{L}_{\mathrm{equiv}}$};
\node[dimlabel, below=3pt of loss, align=center]
    {$\beta$: KL annealing
     \qquad $\lambda_{\mathrm{equiv}}$: equivariance weight};

\draw[->, line width=0.9pt, draw=red!40!black]
    (dec.south) -- (dec.south |- loss.north);

% =====================================================================
%  EQUIVARIANCE REGULARIZATION
% =====================================================================
\node[rectangle, rounded corners=4pt, equiv_col,
      minimum height=1.65cm, minimum width=4.8cm,
      align=center] (equiv) at (10.4, -8.5) {};
\node[seclabel, text=teal!45!black, anchor=north]
    at ([yshift=-4pt]equiv.north) {Equivariance};
\node[font=\footnotesize, align=center] at ([yshift=-0.1cm]equiv.center)
    {$\mathcal{D}\bigl(\mathcal{T}(v); c\bigr)
      \approx
      \mathcal{T}\bigl(\mathcal{D}(v; c)\bigr)$\\[2pt]
     {\scriptsize\itshape decode$\,\circ\,$transform
      $\approx$
      transform$\,\circ\,$decode}};

\draw[dasharrow] (equiv.west) -- (loss.east);

\end{tikzpicture}
}% end resizebox

\caption{
Architecture of L-GTA. The \textbf{Encoder}~$\mathcal{E}$ processes each window through a Bi-LSTM, followed by temporal multi-head self-attention with learned positional embeddings, producing per-timestep Gaussian parameters $(\mu_t, \Sigma_t)$. Latent vectors $v_t$ are sampled via the reparameterization trick. Parametric \textbf{latent-space transformations} $\mathcal{T}_\eta$ are then applied to $v$ and can be chained arbitrarily. The \textbf{Decoder}~$\mathcal{D}$ reconstructs the series from the transformed latent $v'$ and context~$c$ via a Bi-LSTM with a learned skip connection ($z_{\mathrm{proj}}$).
An \textbf{equivariance} term $\mathcal{L}_{\mathrm{equiv}}$ regularizes the decoder so that $\mathcal{D}(\mathcal{T}(v); c) \approx \mathcal{T}(\mathcal{D}(v; c))$.
}
\label{fig:lgta_architecture}
\end{figure*}


\begin{figure*}[t]
\centering
\resizebox{\textwidth}{!}{%
\begin{tikzpicture}[
    >=Stealth,
    % ── colour palette ──
    encoder_col/.style  = {fill=blue!6,   draw=blue!45!black,  line width=0.8pt},
    decoder_col/.style  = {fill=orange!6, draw=orange!50!black, line width=0.8pt},
    latent_col/.style   = {fill=green!6,  draw=green!45!black,  line width=0.8pt},
    transf_col/.style   = {fill=violet!6, draw=violet!45!black, line width=0.8pt},
    io_col/.style       = {fill=gray!6,   draw=gray!55,         line width=0.8pt},
    loss_col/.style     = {fill=red!5,    draw=red!45!black,    line width=0.8pt},
    equiv_col/.style    = {fill=teal!5,   draw=teal!50!black,   line width=0.8pt,
                           dashed, dash pattern=on 3pt off 2pt},
    ctx_col/.style      = {fill=yellow!6, draw=yellow!50!black, line width=0.8pt},
    % ── shapes ──
    bigblock/.style = {
        rectangle, rounded corners=4pt,
        minimum height=3.4cm, minimum width=4.2cm,
        align=center},
    medblock/.style = {
        rectangle, rounded corners=4pt,
        minimum height=3.4cm, minimum width=3.6cm,
        align=center},
    innerblock/.style = {
        rectangle, rounded corners=2pt,
        minimum height=0.6cm, minimum width=3.0cm,
        align=center, font=\footnotesize, line width=0.5pt},
    ioblock/.style = {
        rectangle, rounded corners=3pt,
        minimum height=1.2cm, minimum width=1.8cm,
        align=center, font=\normalsize, io_col},
    prepblock/.style = {
        rectangle, rounded corners=3pt,
        minimum height=1.2cm, minimum width=2.4cm,
        align=center, font=\footnotesize, io_col},
    arrow/.style = {->, line width=1.2pt, draw=gray!60},
    thinarrow/.style = {->, line width=0.8pt, draw=gray!50},
    dasharrow/.style = {->, line width=0.8pt, draw=teal!50!black,
                        dashed, dash pattern=on 3pt off 2pt},
    dimlabel/.style = {font=\footnotesize\itshape, text=gray!65!black},
    seclabel/.style = {font=\normalsize\bfseries},
]

% =====================================================================
%  ROW 1 (top, left to right): Input → Encoder → Latent
% =====================================================================

\node[ioblock] (input) at (0, 0)
    {Input\\[2pt]$z \in \mathbb{R}^{T \times S}$};

\node[bigblock, encoder_col] (enc) at (5.4, 0) {};
\node[seclabel, text=blue!45!black, anchor=north]
    at ([yshift=-4pt]enc.north) {Encoder\; $\mathcal{E}$};
\node[innerblock, encoder_col] (enc_bilstm)
    at ([yshift=0.05cm]enc.center)
    {Bi-LSTM $\to$ LayerNorm};
\node[innerblock, encoder_col, below=0.3cm of enc_bilstm] (enc_attn)
    {Temporal Self-Attention\\[-1pt]
     {\scriptsize pos.\ embed \textbullet{} MHA \textbullet{} FFN}};

\node[medblock, latent_col, minimum height=3.8cm] (lat) at (10.4, 0) {};
\node[seclabel, text=green!38!black, anchor=north]
    at ([yshift=-4pt]lat.north) {Latent Space};
\node[innerblock, latent_col, minimum width=2.4cm]
    (lat_mu) at ([yshift=0.1cm]lat.center)
    {$\mu_t,\;\log\Sigma_t$};
\node[innerblock, latent_col, minimum width=2.4cm,
      below=0.25cm of lat_mu] (lat_sample)
    {Sampling\\[-1pt]{\scriptsize reparam.\ trick}};
\node[dimlabel, below=0.1cm of lat_sample]
    {$v_t \sim \mathcal{N}(\mu_t,\, \Sigma_t)$};

% =====================================================================
%  TURN: Latent Space → Transforms (going down)
% =====================================================================

\node[medblock, transf_col, minimum height=4.2cm, minimum width=3.9cm] (transf) at (10.4, -5.2) {};
\node[seclabel, text=violet!45!black, anchor=north]
    at ([yshift=-4pt]transf.north) {Latent Transforms};
\node[innerblock, transf_col, minimum width=2.4cm]
    (transf_fn) at ([yshift=0.45cm]transf.center)
    {$v' = \mathcal{T}_\eta(v,\;\eta)$};
\node[font=\footnotesize, text=violet!40!black, align=center,
      below=0.25cm of transf_fn] (transf_list)
    {jitter \textbullet{} scaling \textbullet{} mag.\ warp\\
     drift \textbullet{} trend};
\node[dimlabel, align=center, below=0.25cm of transf_list]
    {chainable: $\mathcal{T}_n \!\circ\! \cdots \!\circ\! \mathcal{T}_1$};

% =====================================================================
%  ROW 2 (bottom, right to left): Decoder → Output
% =====================================================================

\node[bigblock, decoder_col, minimum height=4.2cm] (dec) at (5.4, -5.2) {};
\node[seclabel, text=orange!50!black, anchor=north]
    at ([yshift=-4pt]dec.north) {Decoder\; $\mathcal{D}$};
\node[font=\footnotesize\itshape, text=orange!45!black,
      anchor=north] at ([yshift=-0.5cm]dec.north)
    {$+\; z_{\mathrm{proj}}(v')$ skip connection};
\node[innerblock, decoder_col] (dec_bilstm)
    at ([yshift=-0.05cm]dec.center) {Bi-LSTM};
\node[innerblock, decoder_col, below=0.3cm of dec_bilstm] (dec_fc)
    {FC };

\node[ioblock] (output) at (0, -5.2)
    {Output\\[2pt]$\tilde{z} \in \mathbb{R}^{T \times S}$};

% =====================================================================
%  CONTEXT (aligned with encoder and decoder)
% =====================================================================
\node[prepblock, ctx_col, minimum width=2.8cm, minimum height=1.2cm]
    (ctx) at (5.4, -2.4)
    {Context $c$\\[1pt]{\scriptsize trigonometric time features}};

% =====================================================================
%  MAIN FLOW ARROWS
% =====================================================================

% Top row (left → right)
\draw[arrow] (input) -- (enc);
\draw[arrow] (enc)   -- (lat);

% Encoder internal
\draw[thinarrow] (enc_bilstm) -- (enc_attn);

% Turn down (latent → transforms)
\draw[arrow] (lat.south) -- (transf.north);

% Bottom row (right → left)
\draw[arrow] (transf.west) -- (dec.east);
\draw[arrow] (dec) -- (output);

% =====================================================================
%  CONTEXT ARROWS: top of context → encoder bottom, bottom of context → decoder top
% =====================================================================
\draw[->, line width=1.0pt, draw=yellow!55!black]
    (ctx.north) -- (enc.south);

\draw[->, line width=1.0pt, draw=yellow!55!black]
    (ctx.south) -- (dec.north);

% =====================================================================
%  TRAINING LOSS
% =====================================================================
\node[rectangle, rounded corners=4pt, loss_col,
      minimum height=1.1cm,
      align=center, font=\small, inner sep=8pt]
    (loss) at (4.8, -8.5)
    {$\mathcal{L}
       = \mathcal{L}_{\mathrm{recon}}
       + \beta\,\mathcal{L}_{\mathrm{KL}}
       + \lambda_{\mathrm{equiv}}\,\mathcal{L}_{\mathrm{equiv}}$};
\node[dimlabel, below=3pt of loss, align=center]
    {$\beta$: KL annealing
     \qquad $\lambda_{\mathrm{equiv}}$: equivariance weight};

\draw[->, line width=0.9pt, draw=red!40!black]
    (dec.south) -- (dec.south |- loss.north);

% =====================================================================
%  EQUIVARIANCE REGULARIZATION
% =====================================================================
\node[rectangle, rounded corners=4pt, equiv_col,
      minimum height=1.65cm, minimum width=4.8cm,
      align=center] (equiv) at (10.4, -8.5) {};
\node[seclabel, text=teal!45!black, anchor=north]
    at ([yshift=-4pt]equiv.north) {Equivariance};
\node[font=\footnotesize, align=center] at ([yshift=-0.1cm]equiv.center)
    {$\mathcal{D}\bigl(\mathcal{T}(v); c\bigr)
      \approx
      \mathcal{T}\bigl(\mathcal{D}(v; c)\bigr)$\\[2pt]
     {\scriptsize\itshape decode$\,\circ\,$transform
      $\approx$
      transform$\,\circ\,$decode}};

\draw[dasharrow] (equiv.west) -- (loss.east);

\end{tikzpicture}
}% end resizebox

\caption{
Architecture of L-GTA. The \textbf{Encoder}~$\mathcal{E}$ processes each window through a Bi-LSTM, followed by temporal multi-head self-attention with learned positional embeddings, producing per-timestep Gaussian parameters $(\mu_t, \Sigma_t)$. Latent vectors $v_t$ are sampled via the reparameterization trick. Parametric \textbf{latent-space transformations} $\mathcal{T}_\eta$ are then applied to $v$ and can be chained arbitrarily. The \textbf{Decoder}~$\mathcal{D}$ reconstructs the series from the transformed latent $v'$ and context~$c$ via a Bi-LSTM with a learned skip connection ($z_{\mathrm{proj}}$).
An \textbf{equivariance} term $\mathcal{L}_{\mathrm{equiv}}$ regularizes the decoder so that $\mathcal{D}(\mathcal{T}(v); c) \approx \mathcal{T}(\mathcal{D}(v; c))$.
}
\label{fig:lgta_architecture}
\end{figure*}


\begin{figure*}[t]
\centering
\resizebox{\textwidth}{!}{%
\begin{tikzpicture}[
    >=Stealth,
    % ── colour palette ──
    encoder_col/.style  = {fill=blue!6,   draw=blue!45!black,  line width=0.8pt},
    decoder_col/.style  = {fill=orange!6, draw=orange!50!black, line width=0.8pt},
    latent_col/.style   = {fill=green!6,  draw=green!45!black,  line width=0.8pt},
    transf_col/.style   = {fill=violet!6, draw=violet!45!black, line width=0.8pt},
    io_col/.style       = {fill=gray!6,   draw=gray!55,         line width=0.8pt},
    loss_col/.style     = {fill=red!5,    draw=red!45!black,    line width=0.8pt},
    equiv_col/.style    = {fill=teal!5,   draw=teal!50!black,   line width=0.8pt,
                           dashed, dash pattern=on 3pt off 2pt},
    ctx_col/.style      = {fill=yellow!6, draw=yellow!50!black, line width=0.8pt},
    % ── shapes ──
    bigblock/.style = {
        rectangle, rounded corners=4pt,
        minimum height=3.4cm, minimum width=4.2cm,
        align=center},
    medblock/.style = {
        rectangle, rounded corners=4pt,
        minimum height=3.4cm, minimum width=3.6cm,
        align=center},
    innerblock/.style = {
        rectangle, rounded corners=2pt,
        minimum height=0.6cm, minimum width=3.0cm,
        align=center, font=\footnotesize, line width=0.5pt},
    ioblock/.style = {
        rectangle, rounded corners=3pt,
        minimum height=1.2cm, minimum width=1.8cm,
        align=center, font=\normalsize, io_col},
    prepblock/.style = {
        rectangle, rounded corners=3pt,
        minimum height=1.2cm, minimum width=2.4cm,
        align=center, font=\footnotesize, io_col},
    arrow/.style = {->, line width=1.2pt, draw=gray!60},
    thinarrow/.style = {->, line width=0.8pt, draw=gray!50},
    dasharrow/.style = {->, line width=0.8pt, draw=teal!50!black,
                        dashed, dash pattern=on 3pt off 2pt},
    dimlabel/.style = {font=\footnotesize\itshape, text=gray!65!black},
    seclabel/.style = {font=\normalsize\bfseries},
]

% =====================================================================
%  ROW 1 (top, left to right): Input → Encoder → Latent
% =====================================================================

\node[ioblock] (input) at (0, 0)
    {Input\\[2pt]$z \in \mathbb{R}^{T \times S}$};

\node[bigblock, encoder_col] (enc) at (5.4, 0) {};
\node[seclabel, text=blue!45!black, anchor=north]
    at ([yshift=-4pt]enc.north) {Encoder\; $\mathcal{E}$};
\node[innerblock, encoder_col] (enc_bilstm)
    at ([yshift=0.05cm]enc.center)
    {Bi-LSTM $\to$ LayerNorm};
\node[innerblock, encoder_col, below=0.3cm of enc_bilstm] (enc_attn)
    {Temporal Self-Attention\\[-1pt]
     {\scriptsize pos.\ embed \textbullet{} MHA \textbullet{} FFN}};

\node[medblock, latent_col, minimum height=3.8cm] (lat) at (10.4, 0) {};
\node[seclabel, text=green!38!black, anchor=north]
    at ([yshift=-4pt]lat.north) {Latent Space};
\node[innerblock, latent_col, minimum width=2.4cm]
    (lat_mu) at ([yshift=0.1cm]lat.center)
    {$\mu_t,\;\log\Sigma_t$};
\node[innerblock, latent_col, minimum width=2.4cm,
      below=0.25cm of lat_mu] (lat_sample)
    {Sampling\\[-1pt]{\scriptsize reparam.\ trick}};
\node[dimlabel, below=0.1cm of lat_sample]
    {$v_t \sim \mathcal{N}(\mu_t,\, \Sigma_t)$};

% =====================================================================
%  TURN: Latent Space → Transforms (going down)
% =====================================================================

\node[medblock, transf_col, minimum height=4.2cm, minimum width=3.9cm] (transf) at (10.4, -5.2) {};
\node[seclabel, text=violet!45!black, anchor=north]
    at ([yshift=-4pt]transf.north) {Latent Transforms};
\node[innerblock, transf_col, minimum width=2.4cm]
    (transf_fn) at ([yshift=0.45cm]transf.center)
    {$v' = \mathcal{T}_\eta(v,\;\eta)$};
\node[font=\footnotesize, text=violet!40!black, align=center,
      below=0.25cm of transf_fn] (transf_list)
    {jitter \textbullet{} scaling \textbullet{} mag.\ warp\\
     drift \textbullet{} trend};
\node[dimlabel, align=center, below=0.25cm of transf_list]
    {chainable: $\mathcal{T}_n \!\circ\! \cdots \!\circ\! \mathcal{T}_1$};

% =====================================================================
%  ROW 2 (bottom, right to left): Decoder → Output
% =====================================================================

\node[bigblock, decoder_col, minimum height=4.2cm] (dec) at (5.4, -5.2) {};
\node[seclabel, text=orange!50!black, anchor=north]
    at ([yshift=-4pt]dec.north) {Decoder\; $\mathcal{D}$};
\node[font=\footnotesize\itshape, text=orange!45!black,
      anchor=north] at ([yshift=-0.5cm]dec.north)
    {$+\; z_{\mathrm{proj}}(v')$ skip connection};
\node[innerblock, decoder_col] (dec_bilstm)
    at ([yshift=-0.05cm]dec.center) {Bi-LSTM};
\node[innerblock, decoder_col, below=0.3cm of dec_bilstm] (dec_fc)
    {FC };

\node[ioblock] (output) at (0, -5.2)
    {Output\\[2pt]$\tilde{z} \in \mathbb{R}^{T \times S}$};

% =====================================================================
%  CONTEXT (aligned with encoder and decoder)
% =====================================================================
\node[prepblock, ctx_col, minimum width=2.8cm, minimum height=1.2cm]
    (ctx) at (5.4, -2.4)
    {Context $c$\\[1pt]{\scriptsize trigonometric time features}};

% =====================================================================
%  MAIN FLOW ARROWS
% =====================================================================

% Top row (left → right)
\draw[arrow] (input) -- (enc);
\draw[arrow] (enc)   -- (lat);

% Encoder internal
\draw[thinarrow] (enc_bilstm) -- (enc_attn);

% Turn down (latent → transforms)
\draw[arrow] (lat.south) -- (transf.north);

% Bottom row (right → left)
\draw[arrow] (transf.west) -- (dec.east);
\draw[arrow] (dec) -- (output);

% =====================================================================
%  CONTEXT ARROWS: top of context → encoder bottom, bottom of context → decoder top
% =====================================================================
\draw[->, line width=1.0pt, draw=yellow!55!black]
    (ctx.north) -- (enc.south);

\draw[->, line width=1.0pt, draw=yellow!55!black]
    (ctx.south) -- (dec.north);

% =====================================================================
%  TRAINING LOSS
% =====================================================================
\node[rectangle, rounded corners=4pt, loss_col,
      minimum height=1.1cm,
      align=center, font=\small, inner sep=8pt]
    (loss) at (4.8, -8.5)
    {$\mathcal{L}
       = \mathcal{L}_{\mathrm{recon}}
       + \beta\,\mathcal{L}_{\mathrm{KL}}
       + \lambda_{\mathrm{equiv}}\,\mathcal{L}_{\mathrm{equiv}}$};
\node[dimlabel, below=3pt of loss, align=center]
    {$\beta$: KL annealing
     \qquad $\lambda_{\mathrm{equiv}}$: equivariance weight};

\draw[->, line width=0.9pt, draw=red!40!black]
    (dec.south) -- (dec.south |- loss.north);

% =====================================================================
%  EQUIVARIANCE REGULARIZATION
% =====================================================================
\node[rectangle, rounded corners=4pt, equiv_col,
      minimum height=1.65cm, minimum width=4.8cm,
      align=center] (equiv) at (10.4, -8.5) {};
\node[seclabel, text=teal!45!black, anchor=north]
    at ([yshift=-4pt]equiv.north) {Equivariance};
\node[font=\footnotesize, align=center] at ([yshift=-0.1cm]equiv.center)
    {$\mathcal{D}\bigl(\mathcal{T}(v); c\bigr)
      \approx
      \mathcal{T}\bigl(\mathcal{D}(v; c)\bigr)$\\[2pt]
     {\scriptsize\itshape decode$\,\circ\,$transform
      $\approx$
      transform$\,\circ\,$decode}};

\draw[dasharrow] (equiv.west) -- (loss.east);

\end{tikzpicture}
}% end resizebox

\caption{
Architecture of L-GTA. The \textbf{Encoder}~$\mathcal{E}$ processes each window through a Bi-LSTM, followed by temporal multi-head self-attention with learned positional embeddings, producing per-timestep Gaussian parameters $(\mu_t, \Sigma_t)$. Latent vectors $v_t$ are sampled via the reparameterization trick. Parametric \textbf{latent-space transformations} $\mathcal{T}_\eta$ are then applied to $v$ and can be chained arbitrarily. The \textbf{Decoder}~$\mathcal{D}$ reconstructs the series from the transformed latent $v'$ and context~$c$ via a Bi-LSTM with a learned skip connection ($z_{\mathrm{proj}}$).
An \textbf{equivariance} term $\mathcal{L}_{\mathrm{equiv}}$ regularizes the decoder so that $\mathcal{D}(\mathcal{T}(v); c) \approx \mathcal{T}(\mathcal{D}(v; c))$.
}
\label{fig:lgta_architecture}
\end{figure*}


\begin{figure*}[t]
\centering
\resizebox{\textwidth}{!}{%
\begin{tikzpicture}[
    >=Stealth,
    % ── colour palette ──
    encoder_col/.style  = {fill=blue!6,   draw=blue!45!black,  line width=0.8pt},
    decoder_col/.style  = {fill=orange!6, draw=orange!50!black, line width=0.8pt},
    latent_col/.style   = {fill=green!6,  draw=green!45!black,  line width=0.8pt},
    transf_col/.style   = {fill=violet!6, draw=violet!45!black, line width=0.8pt},
    io_col/.style       = {fill=gray!6,   draw=gray!55,         line width=0.8pt},
    loss_col/.style     = {fill=red!5,    draw=red!45!black,    line width=0.8pt},
    equiv_col/.style    = {fill=teal!5,   draw=teal!50!black,   line width=0.8pt,
                           dashed, dash pattern=on 3pt off 2pt},
    ctx_col/.style      = {fill=yellow!6, draw=yellow!50!black, line width=0.8pt},
    % ── shapes ──
    bigblock/.style = {
        rectangle, rounded corners=4pt,
        minimum height=3.4cm, minimum width=4.2cm,
        align=center},
    medblock/.style = {
        rectangle, rounded corners=4pt,
        minimum height=3.4cm, minimum width=3.6cm,
        align=center},
    innerblock/.style = {
        rectangle, rounded corners=2pt,
        minimum height=0.6cm, minimum width=3.0cm,
        align=center, font=\footnotesize, line width=0.5pt},
    ioblock/.style = {
        rectangle, rounded corners=3pt,
        minimum height=1.2cm, minimum width=1.8cm,
        align=center, font=\normalsize, io_col},
    prepblock/.style = {
        rectangle, rounded corners=3pt,
        minimum height=1.2cm, minimum width=2.4cm,
        align=center, font=\footnotesize, io_col},
    arrow/.style = {->, line width=1.2pt, draw=gray!60},
    thinarrow/.style = {->, line width=0.8pt, draw=gray!50},
    dasharrow/.style = {->, line width=0.8pt, draw=teal!50!black,
                        dashed, dash pattern=on 3pt off 2pt},
    dimlabel/.style = {font=\footnotesize\itshape, text=gray!65!black},
    seclabel/.style = {font=\normalsize\bfseries},
]

% =====================================================================
%  ROW 1 (top, left to right): Input → Encoder → Latent
% =====================================================================

\node[ioblock] (input) at (0, 0)
    {Input\\[2pt]$z \in \mathbb{R}^{T \times S}$};

\node[bigblock, encoder_col] (enc) at (5.4, 0) {};
\node[seclabel, text=blue!45!black, anchor=north]
    at ([yshift=-4pt]enc.north) {Encoder\; $\mathcal{E}$};
\node[innerblock, encoder_col] (enc_bilstm)
    at ([yshift=0.05cm]enc.center)
    {Bi-LSTM $\to$ LayerNorm};
\node[innerblock, encoder_col, below=0.3cm of enc_bilstm] (enc_attn)
    {Temporal Self-Attention\\[-1pt]
     {\scriptsize pos.\ embed \textbullet{} MHA \textbullet{} FFN}};

\node[medblock, latent_col, minimum height=3.8cm] (lat) at (10.4, 0) {};
\node[seclabel, text=green!38!black, anchor=north]
    at ([yshift=-4pt]lat.north) {Latent Space};
\node[innerblock, latent_col, minimum width=2.4cm]
    (lat_mu) at ([yshift=0.1cm]lat.center)
    {$\mu_t,\;\log\Sigma_t$};
\node[innerblock, latent_col, minimum width=2.4cm,
      below=0.25cm of lat_mu] (lat_sample)
    {Sampling\\[-1pt]{\scriptsize reparam.\ trick}};
\node[dimlabel, below=0.1cm of lat_sample]
    {$v_t \sim \mathcal{N}(\mu_t,\, \Sigma_t)$};

% =====================================================================
%  TURN: Latent Space → Transforms (going down)
% =====================================================================

\node[medblock, transf_col, minimum height=4.2cm, minimum width=3.9cm] (transf) at (10.4, -5.2) {};
\node[seclabel, text=violet!45!black, anchor=north]
    at ([yshift=-4pt]transf.north) {Latent Transforms};
\node[innerblock, transf_col, minimum width=2.4cm]
    (transf_fn) at ([yshift=0.45cm]transf.center)
    {$v' = \mathcal{T}_\eta(v,\;\eta)$};
\node[font=\footnotesize, text=violet!40!black, align=center,
      below=0.25cm of transf_fn] (transf_list)
    {jitter \textbullet{} scaling \textbullet{} mag.\ warp\\
     drift \textbullet{} trend};
\node[dimlabel, align=center, below=0.25cm of transf_list]
    {chainable: $\mathcal{T}_n \!\circ\! \cdots \!\circ\! \mathcal{T}_1$};

% =====================================================================
%  ROW 2 (bottom, right to left): Decoder → Output
% =====================================================================

\node[bigblock, decoder_col, minimum height=4.2cm] (dec) at (5.4, -5.2) {};
\node[seclabel, text=orange!50!black, anchor=north]
    at ([yshift=-4pt]dec.north) {Decoder\; $\mathcal{D}$};
\node[font=\footnotesize\itshape, text=orange!45!black,
      anchor=north] at ([yshift=-0.5cm]dec.north)
    {$+\; z_{\mathrm{proj}}(v')$ skip connection};
\node[innerblock, decoder_col] (dec_bilstm)
    at ([yshift=-0.05cm]dec.center) {Bi-LSTM};
\node[innerblock, decoder_col, below=0.3cm of dec_bilstm] (dec_fc)
    {FC };

\node[ioblock] (output) at (0, -5.2)
    {Output\\[2pt]$\tilde{z} \in \mathbb{R}^{T \times S}$};

% =====================================================================
%  CONTEXT (aligned with encoder and decoder)
% =====================================================================
\node[prepblock, ctx_col, minimum width=2.8cm, minimum height=1.2cm]
    (ctx) at (5.4, -2.4)
    {Context $c$\\[1pt]{\scriptsize trigonometric time features}};

% =====================================================================
%  MAIN FLOW ARROWS
% =====================================================================

% Top row (left → right)
\draw[arrow] (input) -- (enc);
\draw[arrow] (enc)   -- (lat);

% Encoder internal
\draw[thinarrow] (enc_bilstm) -- (enc_attn);

% Turn down (latent → transforms)
\draw[arrow] (lat.south) -- (transf.north);

% Bottom row (right → left)
\draw[arrow] (transf.west) -- (dec.east);
\draw[arrow] (dec) -- (output);

% =====================================================================
%  CONTEXT ARROWS: top of context → encoder bottom, bottom of context → decoder top
% =====================================================================
\draw[->, line width=1.0pt, draw=yellow!55!black]
    (ctx.north) -- (enc.south);

\draw[->, line width=1.0pt, draw=yellow!55!black]
    (ctx.south) -- (dec.north);

% =====================================================================
%  TRAINING LOSS
% =====================================================================
\node[rectangle, rounded corners=4pt, loss_col,
      minimum height=1.1cm,
      align=center, font=\small, inner sep=8pt]
    (loss) at (4.8, -8.5)
    {$\mathcal{L}
       = \mathcal{L}_{\mathrm{recon}}
       + \beta\,\mathcal{L}_{\mathrm{KL}}
       + \lambda_{\mathrm{equiv}}\,\mathcal{L}_{\mathrm{equiv}}$};
\node[dimlabel, below=3pt of loss, align=center]
    {$\beta$: KL annealing
     \qquad $\lambda_{\mathrm{equiv}}$: equivariance weight};

\draw[->, line width=0.9pt, draw=red!40!black]
    (dec.south) -- (dec.south |- loss.north);

% =====================================================================
%  EQUIVARIANCE REGULARIZATION
% =====================================================================
\node[rectangle, rounded corners=4pt, equiv_col,
      minimum height=1.65cm, minimum width=4.8cm,
      align=center] (equiv) at (10.4, -8.5) {};
\node[seclabel, text=teal!45!black, anchor=north]
    at ([yshift=-4pt]equiv.north) {Equivariance};
\node[font=\footnotesize, align=center] at ([yshift=-0.1cm]equiv.center)
    {$\mathcal{D}\bigl(\mathcal{T}(v); c\bigr)
      \approx
      \mathcal{T}\bigl(\mathcal{D}(v; c)\bigr)$\\[2pt]
     {\scriptsize\itshape decode$\,\circ\,$transform
      $\approx$
      transform$\,\circ\,$decode}};

\draw[dasharrow] (equiv.west) -- (loss.east);

\end{tikzpicture}
}% end resizebox

\caption{
Architecture of L-GTA. The \textbf{Encoder}~$\mathcal{E}$ processes each window through a Bi-LSTM, followed by temporal multi-head self-attention with learned positional embeddings, producing per-timestep Gaussian parameters $(\mu_t, \Sigma_t)$. Latent vectors $v_t$ are sampled via the reparameterization trick. Parametric \textbf{latent-space transformations} $\mathcal{T}_\eta$ are then applied to $v$ and can be chained arbitrarily. The \textbf{Decoder}~$\mathcal{D}$ reconstructs the series from the transformed latent $v'$ and context~$c$ via a Bi-LSTM with a learned skip connection ($z_{\mathrm{proj}}$).
An \textbf{equivariance} term $\mathcal{L}_{\mathrm{equiv}}$ regularizes the decoder so that $\mathcal{D}(\mathcal{T}(v); c) \approx \mathcal{T}(\mathcal{D}(v; c))$.
}
\label{fig:lgta_architecture}
\end{figure*}
